\chapter{Approximate inference with noisy-MCMC}
\label{chap:approximate}


Current approaches:
\begin{itemize}
    \item Function emulation: \cite{adame2025, carvalho2021}
\end{itemize}

\comLM{\cite{han2023} \cite{carvalho2021} and references therein }

\comLM{Need to carefully read, cite and discuss \cite{park2020}. Do their results hold for SCAM, for instance?}


Our contributions in this area:
\begin{itemize}
    \item Do MaL estimation and talk about GP and SCAM;
    \item Importance sampling to simplify and improve estimation of the NPP normalising constant (including control variates);
    \item Showcase on moderate size examples (including NIG linear regression as a baseline). Show that it works and then show how to break it    
\end{itemize}


\section{Emulating the normalizing constant}

In many applications involving normalized power priors, the normalizing constant

\begin{equation}
    c_0(\eta)
    =
    \int_{\Theta}
    L(D_0 \mid \theta)^\eta
    \pi_0(\theta)
    \dd \theta
\end{equation}

cannot be evaluated analytically.
As a consequence, direct evaluation of the posterior density becomes infeasible.
A common strategy is therefore to replace $c_0(\eta)$ by a surrogate approximation constructed from evaluations at a finite collection of design points.

More precisely, suppose we evaluate the log-normalizing constant

\begin{equation}
    g(\eta)
    =
    \log c_0(\eta)
\end{equation}

at points $\eta_1, \ldots, \eta_m \in [0,1]$.
The evaluations may be obtained analytically, via Monte Carlo estimators, thermodynamic integration, or path sampling methods.
We then construct an emulator

\begin{equation}
    \widehat g(\eta)
\end{equation}

intended to approximate $g(\eta)$ over the entire interval.
This yields the approximation

\begin{equation}
    c_0(\eta)
    \approx
    \exp\{\widehat g(\eta)\}.
\end{equation}

The resulting approximation can then be plugged into the posterior density, producing an approximate MCMC scheme whose computational cost is substantially reduced after the emulator is trained.

From a statistical perspective, this transforms the problem of evaluating an intractable normalizing constant into a nonparametric regression problem.
The quality of inference then depends on the approximation properties of the emulator and on how uncertainty propagates through the posterior distribution.

\subsection{Consequences}

Although emulation methods can dramatically reduce computational cost, replacing the true normalizing constant by a deterministic approximation may substantially alter the resulting posterior distribution.

Suppose the true posterior density is

\begin{equation}
    \pi(\theta, \eta \mid D, D_0)
    \propto
    \frac{
        L(D \mid \theta)
        L(D_0 \mid \theta)^\eta
        \pi_0(\theta)
        \pi_0(\eta)
    }{
        c_0(\eta)
    }.
\end{equation}

If $c_0(\eta)$ is replaced by an approximation $\widehat c_0(\eta)$, the resulting Markov chain targets instead

\begin{equation}
    \widetilde \pi(\theta, \eta \mid D, D_0)
    \propto
    \frac{
        L(D \mid \theta)
        L(D_0 \mid \theta)^\eta
        \pi_0(\theta)
        \pi_0(\eta)
    }{
        \widehat c_0(\eta)
    }.
\end{equation}

Consequently, approximation errors directly perturb the invariant distribution of the Markov chain.
Even small systematic biases in $\widehat c_0(\eta)$ may produce substantial distortions in posterior summaries, credible intervals, and Bayes factors.

An important practical issue is that many emulation approaches treat $\widehat c_0(\eta)$ as deterministic after fitting.
In this case, uncertainty associated with the emulator is ignored during posterior inference.
This may lead to overconfident posterior distributions and unreliable uncertainty quantification.

The problem becomes particularly severe when the approximation error varies non-uniformly across the parameter space.
For instance, local interpolation errors near regions of high posterior mass may induce systematic bias in the marginal posterior distribution of $\eta$.

Moreover, unconstrained emulators may violate structural properties known to hold for the true function.
For normalized power priors, the map

\begin{equation}
    g(\eta)
    =
    \log c_0(\eta)
\end{equation}

typically satisfies smoothness and monotonicity properties inherited from the likelihood.
Approximations that ignore such information may generate unrealistic artifacts, including oscillations and locally inconsistent curvature.

These issues motivate the incorporation of structural constraints into the emulation procedure.

\subsection{Introducing shape-constraints}

One strategy for improving the stability of the emulator is to incorporate qualitative information about the target function directly into the model.

Suppose

\begin{equation}
    g(\eta)
    =
    \log c_0(\eta).
\end{equation}

Under regularity conditions, derivatives of $g$ can be expressed in terms of moments under the power posterior distribution.
For instance,

\begin{equation}
    \dv{}{\,\eta}
    \log c_0(\eta)
    =
    \bE_{\pi_\eta}
    \qty[
        \log L(D_0 \mid \theta)
    ].
\end{equation}

Therefore, information about monotonicity and curvature may be available analytically or estimated from Monte Carlo samples.

A flexible approach consists of modelling $g$ as a Gaussian process.
If

\begin{equation}
    g
    \sim
    \mathcal{GP}(m, k),
\end{equation}

then derivatives of $g$ are also Gaussian processes because differentiation is a linear operator.
In particular,

\begin{equation}
    g'
    \sim
    \mathcal{GP}
    \qty(
        \dv{}{\,\eta} m(\eta),
        \dv{}{\,\eta}\dv{}{\,\eta'} k(\eta, \eta')
    ).
\end{equation}

This allows derivative information and shape restrictions to be incorporated directly into the posterior distribution of the Gaussian process.
Conditioning on derivative observations or derivative sign constraints can substantially reduce unrealistic oscillatory behaviour.
We refer to \citeonline{adame2025} for recent developments in shape-constrained emulation through Gaussian processes.
\comEA{Insert more references}

An alternative strategy is based on shape-constrained additive models (SCAMs).
These models extend generalized additive models (GAMs) by imposing constraints directly on the spline basis coefficients.

Suppose

\begin{equation}
    g(\eta)
    =
    \beta_0
    +
    \sum_{j=1}^p
    f_j(\eta),
\end{equation}

where the functions $f_j$ are represented through spline expansions.
Shape constraints such as monotonicity or convexity can then be enforced by restricting the spline coefficients during estimation.

In practice, SCAMs are fitted using penalized likelihood methods analogous to standard GAM estimation procedures.
The main difference is that the optimization is performed under linear inequality constraints on the spline coefficients.

Compared with unconstrained smoothers, shape-constrained models may provide more stable extrapolation behaviour and improved uncertainty quantification in regions with sparse design points.
These properties are particularly appealing when approximating normalizing constants, since small local errors may propagate nonlinearly into the posterior distribution.

\section{Current approaches for the NPP}
\comEA{Discuss \cite{han2023} and \cite{carvalho2021} - and any other state-of-art}

\section{Doubly Intractable Settings}

\comEA{Discuss \cite{park2020}}

Doubly intractable distributions arise in Bayesian inference when the likelihood takes the form
\begin{equation}
\pi(\theta \mid x) \propto p(\theta)\frac{h(x \mid \theta)}{Z(\theta)},
\end{equation}
where the normalizing constant $Z(\theta)$ depends on the parameter $\theta$ and cannot be evaluated pointwise, as discussed in detail on seciton \ref{sec:dbl-int}.
This setting is common in models such as exponential random graph models, Markov point processes, and, more generally, in likelihoods derived from Gibbs measures.
In the context of Normalized Power Priors (NPP), this issue naturally appears due to the presence of parameter-dependent normalizing constants, making standard Bayesian computation challenging.

A number of Markov chain Monte Carlo (MCMC) methods have been proposed to address this problem, including auxiliary variable approaches such as the exchange algorithm and the double Metropolis--Hastings (DMH) algorithm. While these methods can target the exact posterior distribution, they typically require simulating auxiliary variables at each iteration, which becomes computationally prohibitive in high-dimensional or complex models.

An alternative strategy is proposed by \citeonline{park2020}, who introduce a function emulation approach to approximate the intractable normalizing constant. Their method replaces repeated Monte Carlo estimation of $Z(\theta)$ with a two-stage surrogate model.
First, importance sampling is used to obtain noisy estimates of $Z(\theta)$ at a finite set of parameter values.
Then, a Gaussian process (GP) is fitted to the log-normalizing function, providing a fast interpolator that can be evaluated within an MCMC algorithm.
This leads to two variants: one that emulates only the normalizing function (NormEm), and another that emulates the entire likelihood (LikEm).

The main advantage of this approach is computational efficiency.
By shifting the costly Monte Carlo approximations to a precomputation stage, the resulting MCMC algorithm avoids repeated simulation of auxiliary variables, yielding substantial speed-ups compared to exact or pseudo-marginal methods.
Moreover, the approach is naturally parallelizable and extends beyond doubly intractable settings to more general problems involving expensive likelihood evaluations.

However, this efficiency comes at the cost of asymptotic exactness.
The resulting Markov chain targets an approximate posterior distribution, with an error that depends on both the number of importance samples and the number of design points used to train the Gaussian process.
While theoretical bounds on the total variation distance are provided, they do not yield explicit rates, making it difficult to calibrate the approximation in practice.

Several additional limitations remain. First, the method relies on importance sampling estimates, which can exhibit high variance in moderate to high-dimensional settings, potentially degrading the quality of the surrogate.
Second, the Gaussian process approximation assumes a smooth structure for the log-normalizing function and typically ignores heteroskedasticity induced by Monte Carlo noise.
Third, the selection of design points (particles) is critical but performed using heuristic procedures such as short runs of DMH or approximate Bayesian computation, without a principled adaptive strategy.
Finally, the approach does not scale well with the dimension of the parameter space due to the cubic complexity of Gaussian process regression and the increasing difficulty of covering the parameter space.

These limitations suggest several directions for future research.
In particular, combining surrogate-based approximations with exact correction mechanisms, such as pseudo-marginal or delayed-acceptance schemes, could potentially recover asymptotic exactness while retaining computational gains.
Additionally, replacing importance sampling with more robust estimators, and incorporating adaptive or local surrogate models, may improve both stability and scalability.
Such developments are particularly relevant in the context of NPP models, where efficient and reliable inference remains an open challenge.